\section{Review Python}

\begin{frame}[fragile]{Types and Variables}
    Python is \textbf{dynamically typed} --- you can assign any type without declaration.

    \begin{lstlisting}
x    = 42              # int
y    = 3.14            # float
s    = "Hello, AIP2!" # str
lst  = [1, 2, 3, 4, 5] # list
flag = True            # bool

print(type(x), type(y), type(s))
    \end{lstlisting}

    \begin{itemize}
        \item \code{type(obj)} returns the runtime type of any object.
        \item All Python values are objects --- including \code{int}, \code{bool}, \code{str}.
    \end{itemize}
\end{frame}

\begin{frame}[fragile]{Iterables}
    \begin{columns}[t]
    \column{0.5\textwidth}
    \begin{block}{Built-in collections}
    \begin{lstlisting}
my_list  = [1, 2, 3]       # mutable
my_tuple = (1, 2, 3)       # immutable
my_set   = {1, 2, 3}       # unique
my_dict  = {"a": 1, "b": 2}
    \end{lstlisting}
    \end{block}

    \column{0.5\textwidth}
    \begin{block}{Iteration}
    \begin{lstlisting}
for item in my_list:
    print(item)

for k, v in my_dict.items():
    print(k, v)

# slicing
print(my_list[1:3])  # [2, 3]
    \end{lstlisting}
    \end{block}
    \end{columns}

    \begin{block}{Generators --- lazy evaluation}
    \begin{lstlisting}
def count_up_to(n):
    count = 1
    while count <= n:
        yield count
        count += 1

for number in count_up_to(5):
    print(number)   # 1 2 3 4 5
    \end{lstlisting}
    \end{block}
\end{frame}

\begin{frame}[fragile]{Functions}
    Defined with \code{def}; \code{return} sends a value back.

    \begin{lstlisting}
def sign(x):
    if x > 0:
        return "positive"
    elif x < 0:
        return "negative"
    else:
        return "zero"

for num in [10, -5, 0]:
    print(f"{num} is {sign(num)}")
    \end{lstlisting}

    \begin{itemize}
        \item Use f-strings (\code{f"\{expr\}"}) for readable formatting.
        \item Functions are first-class objects --- passable as arguments.
    \end{itemize}
\end{frame}

\begin{frame}[fragile]{Python Libraries}
    \begin{lstlisting}
import math
print(math.sqrt(16))   # 4.0
print(math.pi)         # 3.14159...

import numpy as np
print(np.array([1, 2, 3]))  # [1 2 3]
print(np.pi)                # 3.14159...
    \end{lstlisting}

    \vspace{0.5em}
    Key libraries for this course:
    \begin{itemize}
        \item \code{numpy}         --- numerical arrays, linear algebra
        \item \code{matplotlib}    --- data visualisation
        \item \code{torch}         --- deep learning (from Lab 3 onward)
        \item \code{line-profiler} --- line-by-line profiling
    \end{itemize}
\end{frame}
